\chapter{Evaluation of LLMs and LLM-based Agents}
\label{chap:evaluating-agents}

The performance measure of Section~\ref{subsec:rationality} drives action selection during the run of a classical agent, since the decide step selects the action expected to maximize it. An LLM-based agent has no such measure active during the run. It is not built for one task, and its objective is stated in the prompt as natural language.

Whether the behavior obtained matches the behavior intended is decided only after the run, by an evaluation setup that defines both what to check and how to check it. That criterion is applied to the text an LLM-based agent produces as output, but also to the trajectory one or more LLM-based agents may produce while pursuing their goal. This chapter reviews how that criterion is constructed, first for the language model as the core agent component, then for the LLM-based agent built on top of it. It also reviews the additional evaluation problems that construction introduces.

\section{Text Matching and Benchmarks}
\label{sec:text-matching-benchmarks}

\subsection{Text Matching}
\label{sec:text-matching}

Scoring a language model's output against a predetermined reference text is the earliest form of automatic evaluation for generated text. This is a natural language generation problem, not one specific to agents or large language models: any system that produces text needs a way to be scored against a reference. BLEU and ROUGE are the traditional metrics for this problem. BLEU compares a generated translation to one or more reference translations through a weighted average of n-gram precision, matched phrase by phrase against the reference \cite{papineniBleuMethodAutomatic2002}. ROUGE applies the same surface-level comparison to summarization, counting overlapping n-grams, word sequences, or word pairs between a generated summary and human-written reference summaries \cite{linROUGEPackageAutomatic2004}. Both metrics reduce evaluation to how closely the generated text matches a reference text, and this reference-based comparison was shown to correlate poorly with human judgment, particularly on tasks that require creativity and diversity \cite{liuGEvalNLGEvaluation2023}.

\subsection{Benchmarks as Capability Exams}
\label{sec:benchmarks-capability-exams}

A benchmark of a large language model is a set of tasks chosen to represent a capability, together with a standardized procedure for running every model on those tasks and a metric for scoring each one. The result is packaged into a leaderboard so that models can be compared under the same protocol \cite{liangHolisticEvaluationLanguage2023}. Running a model on a task requires an \emph{adaptation} procedure. Adaptation is the method that turns a general-purpose language model into a solver for that one task, for example a specific prompting strategy \cite{liangHolisticEvaluationLanguage2023}. The comparison across models is meaningful only if every model is evaluated on the same scenarios and adapted to each task by the same procedure; otherwise, a difference in score can come from the adaptation method rather than from the model itself \cite{liangHolisticEvaluationLanguage2023}. A benchmark may score an output with text-matching metrics such as BLEU and ROUGE, or apply a different metric when the task requires checking something other than how closely the output resembles a reference string.

MMLU (Massive Multitask Language Understanding) is a benchmark of this kind, on a task where the answer is a single choice. It is a multiple-choice test built from 57 subjects across the humanities, sciences, and other domains, and each item is scored by classification accuracy on the selected answer choice \cite{hendrycksMeasuringMassiveMultitask2021}. The correct answer is one of a predetermined set of options, so the model is scored by whether it selects the option matching the key, not by any comparison of generated text against a reference string.

BFCL (Berkeley Function Calling Leaderboard), another benchmark, illustrates the same point as MMLU for a task where the answer is not text at all: the metric fits the kind of answer the task produces. It evaluates a model's ability to produce function calls and scores them through abstract-syntax-tree sub-string matching, used as a proxy for executing the call: a function call is correct if the function name matches exactly and every parameter value falls within its allowed range \cite{patilBerkeleyFunctionCalling2025}. A function call is a structured object, not a string to be compared word by word, so the benchmark checks its syntactic structure and argument values instead of applying a text-matching metric to it. MMLU and BFCL score two different kinds of output, classification accuracy on a predetermined answer choice and structural correctness on a generated call, and in both cases the metric is chosen to fit what the task is asking.

This benchmark model presupposes headroom: the leaderboard has to separate models by capability, which requires that the tasks not already be solved by every model under test.

\section{Benchmark Saturation}
\label{sec:benchmark-saturation}

The benchmark Humanity's Last Exam is introduced on the observation that the headroom MMLU relies on has closed: language models have now exceeded 90\% accuracy on popular benchmarks such as MMLU, which limits how much such a benchmark can still measure about state-of-the-art capability \cite{phanHumanitysLastExam2025}. Once most models under test score near the ceiling, the benchmark no longer separates them by capability, and the leaderboard collapses toward a narrow band of scores that no longer reflects differences in what the models can do.

Agent benchmarks report the same closing of headroom: top-performing agents have reached scores exceeding 90\% on existing benchmarks, which are becoming saturated and losing their ability to differentiate state-of-the-art agents \cite{guoSelfEvolvingBenchmarksSynthesizing2026}. Saturation is therefore a general property of predetermined benchmarks as models improve.

Benchmarks presuppose headroom: they separate models by capability only if the tasks are not already solved by every model under test. Headroom can close through mechanisms other than accumulated capability gains. Evaluation data can leak into training corpora, and fine-tuning on it inflates a model's score without a matching gain in capability \cite{akhtarWhenAIBenchmarks2026}. Benchmarks therefore need constant revision, proposing tasks difficult enough to preserve headroom \cite{akhtarWhenAIBenchmarks2026}.

Saturation is one way a predetermined evaluation stops discriminating between models. A second way is for the evaluation to be poorly matched to the domain it is applied to. One example is Singhal et al.'s work on medical question answering, where existing benchmarks are limited to classification accuracy or automated natural language generation metrics such as BLEU, and do not support the detailed analysis that real-world clinical use requires \cite{singhalLargeLanguageModels2023}. In place of such generic metrics, evaluation of clinical answers is carried out along axes specific to the domain, including factuality, comprehension, reasoning, possible harm, and bias \cite{singhalLargeLanguageModels2023}. A generic accuracy score or a generic text-generation metric does not capture what makes a clinical answer acceptable, so the metric is redefined around the domain instead of reused from a general-purpose benchmark.

A domain-tailored metric such as this one is still a single predetermined metric, applied to an output that is checked against a predetermined notion of what counts as acceptable. This is sufficient as long as the task itself has a comparison point of that kind. It stops being sufficient once the task is open-ended and has no predetermined answer to compare the output against, so that no single metric, however well matched to the domain, is left with anything to compare it to.

\section{Holistic Evaluation}
\label{sec:holistic-evaluation}

An open-ended or free-form generation task has no predetermined answer to compare its output against, so neither a predetermined-answer benchmark (Section~\ref{sec:benchmarks-capability-exams}) nor reference-based text matching (Section~\ref{sec:text-matching}) applies to it. The same limit applies to retrieval-augmented generation (RAG). RAG systems are commonly evaluated as extractive question answering, against datasets built around short extractive answers, but this does not represent how a RAG system is actually used, since it generates a synthesized answer rather than returning a predetermined span from the retrieved text \cite{esRagasAutomatedEvaluation2025}. As with open-ended generation, the output cannot be checked against a predetermined reference, so a different evaluation approach is needed.

The Holistic Evaluation of Language Models (HELM) responds to this by evaluating a model along several dimensions at once, rather than reducing it to a single score \cite{liangHolisticEvaluationLanguage2023}. Table~\ref{tab:helm-metrics} lists the seven metrics that HELM records for each scenario, with the property that each one measures \cite{liangHolisticEvaluationLanguage2023}.

\begin{table}[ht]
    \centering
    \renewcommand{\arraystretch}{1.2}
    \begin{tabularx}{0.85\textwidth}{@{}>{\raggedright\arraybackslash}p{2.2cm}X@{}}
        \toprule
        Metric & What it measures \\
        \midrule
        Accuracy & Task-specific correctness of the output, such as exact match or the F1-score against the expected response. \\
        Calibration & How closely the model's predicted confidence matches its actual rate of being correct. \\
        Robustness & Worst-case accuracy under small, meaning-preserving input perturbations, such as typos or paraphrases. \\
        Fairness & Disparity in accuracy across social groups, measured through counterfactual perturbations of group-related terms. \\
        Bias & Asymmetry in the language a model generates, such as skewed demographic representation or stereotyped associations. \\
        Toxicity & Rate of offensive, hateful, or otherwise harmful content in model generations, detected by an automated classifier. \\
        Efficiency & Computational cost of training and inference, including energy use, carbon emissions, and runtime. \\
        \bottomrule
    \end{tabularx}
    \caption[The seven metrics that HELM records for each scenario]{The seven metrics that HELM records for each scenario \cite{liangHolisticEvaluationLanguage2023}.}
    \label{tab:helm-metrics}
\end{table}

Holistic evaluation is defined around three elements, each instantiated concretely in HELM's own design \cite{liangHolisticEvaluationLanguage2023}:

\begin{itemize}
    \item \emph{Broad coverage.} A wide range of scenarios is included, while what is left out is explicitly acknowledged. A scenario is a triple of task, domain, and language that determines a concrete use case to evaluate against, for example summarization of news articles in English \cite{liangHolisticEvaluationLanguage2023}. HELM realizes broad coverage by defining a taxonomy of scenarios and metrics, narrowed down to a core set of 16 scenarios.
    \item \emph{Multi-metric measurement.} Several metrics are measured for each scenario, rather than one. HELM records all 7 metrics from Table~\ref{tab:helm-metrics} for each of the 16 core scenarios, so that metrics other than accuracy are not left out of the evaluation.
    \item \emph{Standardization.} HELM adapts every model to each of the 16 core scenarios through the same few-shot prompts.
\end{itemize}

HELM adds a fourth, separate design choice on top of the core scenario and metric grid: \emph{targeted evaluations}, additional scenarios that isolate one skill or risk at a time, covering linguistic understanding, world and commonsense knowledge, reasoning, memorization and copyright, disinformation generation, bias, and toxicity generation \cite{liangHolisticEvaluationLanguage2023}.


\section{LLM-as-a-Judge}
\label{sec:llm-as-judge}

Each of HELM's 7 metrics from Section~\ref{sec:holistic-evaluation} is computed against a reference the scenario determines in advance. HELM's own evaluation includes two scenarios with no such reference: judging whether a generated headline or social media post supports a given thesis. None of the 7 metrics can be computed for these, so HELM falls back to human annotators recruited on Amazon Mechanical Turk, with each output rated by three workers, because no established automated evaluation exists for disinformation \cite{liangHolisticEvaluationLanguage2023}. Scoring an open-ended output therefore requires an \emph{oracle} in place of a predetermined reference string: a procedure that returns a verdict, given a candidate output, on whether it is a good answer. HELM uses a human rater for this case; a human rater is one such oracle, but human rating at scale is costly and could lead to poor repeatability. Using a language model as a reference-free evaluator has since been proposed by a growing body of work as a scalable alternative to a human rater \cite{liuGEvalNLGEvaluation2023}: instead of a human, a language model reads the candidate output and produces a verdict on its quality, standing in for the human rater or for a task-specific score function. This approach is named LLM-as-a-judge, and the model performing the grading is the judge.

\subsection{Grading Modes and the Reliability Question}
\label{sec:llm-as-judge-reliability}

Zheng et al. define three ways the judge can be used \cite{zhengJudgingLLMasaJudgeMTBench2023}:

\begin{itemize}
    \item \emph{Pairwise comparison}: the judge reads two candidate responses to the same prompt and states which one is better, or declares a tie.
    \item \emph{Single-answer grading}: unlike in pairwise comparison, the judge reads one response, with no other response to compare it against, and assigns it a score on a predetermined scale.
    \item \emph{Reference-guided grading}: the judge reads the candidate response together with a reference answer or a set of grading instructions, and scores the candidate against that reference.
\end{itemize}

Zheng et al. evaluate GPT-4 as a judge on two benchmarks built for this purpose: MT-Bench, a set of 80 multi-turn, open-ended questions spanning categories such as writing, reasoning, and coding, and Chatbot Arena, a crowdsourced platform on which users pose the same prompt to two anonymized models and vote for the response they prefer \cite{zhengJudgingLLMasaJudgeMTBench2023}. On these two benchmarks, GPT-4 agreed with human preference more than 80\% of the time, reaching 85\% on MT-Bench. The authors report this rate as matching, and on MT-Bench slightly exceeding, the 81\% agreement measured between two human annotators \cite{zhengJudgingLLMasaJudgeMTBench2023}. The result reads as evidence that a judge model can stand in for a human rater with little loss of reliability.

Reliability also depends on whether a judge's verdict is stable and free of bias. Shi et al. test this by evaluating fifteen LLM judges across MT-Bench and DevBench, a benchmark of software development tasks, over more than 150,000 evaluation instances \cite{shiJudgingJudgesSystematic2025}. Position bias is the tendency to favor a response because of its position in the comparison, such as being listed first, rather than because of its content. They report that position bias is both judge-dependent and task-dependent: the same judge model does not show the same degree of it across tasks, and different judge models rank differently on the same task \cite{shiJudgingJudgesSystematic2025}. Comparing judges against each other, rather than against human raters, they also find that judges with similar aggregate consistency scores frequently disagree on individual instances, with disagreement concentrated on cases where the candidate responses being compared are close in quality \cite{shiJudgingJudgesSystematic2025}. A bias or consistency figure reported for one judge on one benchmark therefore does not transfer to a different judge or a different benchmark without being measured again.

\subsection{G-Eval: A Framework for Building LLM-based Metrics}
\label{sec:llm-as-judge-geval}

Liu et al. propose G-Eval, not a single metric but a framework for building an LLM-based metric for a given evaluation criterion \cite{liuGEvalNLGEvaluation2023}. A G-Eval metric is assembled from a task introduction, a natural language instruction that states the task being evaluated, and a definition of the evaluation criterion, for example coherence in text summarization. The judge model expands these into its own chain-of-thought of evaluation steps before it is shown the candidate output and asked to score it \cite{liuGEvalNLGEvaluation2023}.

Reading a single score off the judge's output is unreliable in two ways that the authors report. First, one score value tends to dominate the distribution of outputs, which lowers the score's variance. Second, language models tend to return integer scores even when the prompt asks for decimal values, which produces ties that mask real differences between outputs \cite{liuGEvalNLGEvaluation2023}. G-Eval addresses both by reading the judge's probability distribution over the next token, at the position where the judge outputs its score. This is the same distribution over the vocabulary described in Section~\ref{subsec:unembedding}. From this distribution, only the probabilities of the tokens that denote a valid score value, such as the digits 1 to 5, are kept. The final score is the sum of each valid score value multiplied by its probability, rather than the single highest-probability token \cite{liuGEvalNLGEvaluation2023}. This turns a discrete, often-tied score into a continuous one, and is reported to raise the correlation with human judgment over the discrete version on the SummEval benchmark \cite{liuGEvalNLGEvaluation2023}.

With GPT-4 as the judge model, G-Eval reached an average Spearman correlation of 0.514 with human judgment on text summarization, averaged over four rated dimensions. This is a moderate correlation in absolute terms, but it is higher than the correlations reported for reference-based metrics such as ROUGE (0.192), embedding-based metrics such as BERTScore (0.225), and other LLM-based evaluators such as GPTScore (0.417) \cite{liuGEvalNLGEvaluation2023}. The authors report a systematic problem alongside this result: G-Eval assigns higher scores to GPT-3.5-generated summaries than to human-written ones, even in cases where human judges prefer the human-written summary, and attribute this to the judge favoring text similar to what it would generate itself \cite{liuGEvalNLGEvaluation2023}. Building a more accurate scoring function does not remove this kind of systematic deviation from human judgment; the next section places it within a broader taxonomy of such deviations.

\subsection{Systematic Biases in LLM-Judge Verdicts}
\label{sec:llm-as-judge-biases}

The literature has organized systematic deviations from human judgment into a taxonomy of biases. Ye et al. identified twelve named bias types, each a way the judge's verdict can shift for reasons unrelated to the quality of the response being graded \cite{yeJusticePrejudiceQuantifying2024a}. Table~\ref{tab:llm-judge-biases} lists all twelve, with a description of each.

\begin{table}[ht]
    \centering
    \renewcommand{\arraystretch}{1.2}
    \begin{tabularx}{0.9\textwidth}{@{}>{\raggedright\arraybackslash}p{3.4cm}X@{}}
        \toprule
        Bias & Description \\
        \midrule
        Position bias & Favoring a response for the position it occupies in a pairwise comparison, rather than its content. \\
        Verbosity bias & Favoring a longer response independently of its quality. \\
        Compassion-fade bias & Judging the same content differently depending on whether the responses are attributed to named models. \\
        Bandwagon bias & Favoring a response after being told it is the majority or popular choice. \\
        Distraction bias & Being swayed by content in the response that is irrelevant to the criterion being graded. \\
        Fallacy-oversight bias & Failing to penalize a response that contains a logical fallacy. \\
        Authority bias & Favoring a response that cites sources or authorities, regardless of whether the citation supports the claim. \\
        Sentiment bias & Favoring a response for its emotional tone rather than its content. \\
        Diversity bias & Judging content differently depending on demographic or cultural markers present in it. \\
        Chain-of-thought bias & The judge's own evaluation accuracy changes depending on whether it is prompted to reason step by step before returning a verdict. \\
        Self-enhancement bias & Favoring a response the judge model itself, or a model similar to it, is more likely to have produced. \\
        Refinement-aware bias & Judging a response differently once it is known to be a revision of an earlier one. \\
        \bottomrule
    \end{tabularx}
    \caption[The twelve bias types identified in LLM-as-a-judge verdicts]{The twelve bias types Ye et al. identify in LLM-as-a-judge verdicts \cite{yeJusticePrejudiceQuantifying2024a}.}
    \label{tab:llm-judge-biases}
\end{table}

Two of these have been studied as specific cases on their own.
\begin{itemize}
	\item Position bias: Shi et al. found that its strength depends on the quality gap between the two responses compared: the closer the two responses are in quality, the more the judge's verdict is affected by which position each response occupies \cite{shiJudgingJudgesSystematic2025}.
	\item Self-enhancement bias, referred to by Wataoka et al. as self-preference: they reported that GPT-4, used as a judge, exhibits a significant degree of self-preference, and attributed this to language models favoring text that is more familiar to them, including text that the model itself, or a model similar to it, is more likely to have produced \cite{wataokaSelfPreferenceBiasLLMasaJudge2025}.
\end{itemize}

Relative to the reference-based metrics of Section~\ref{sec:text-matching-benchmarks}, an LLM judge can grade an output for which no predetermined reference exists, since the judge substitutes for the reference rather than requiring one. This is what makes it applicable to open-ended generation. The cost of this flexibility is exposure to the biases described above. A reference-based metric scores one output at a time against its predetermined reference string, so the position of a response in a comparison cannot affect its score. It also does not generate text itself, so it cannot favor a response for resembling its own output, unlike a judge model. An LLM judge reads the full context of the comparison, so it is exposed to signals that the grading task was not meant to depend on. The rate of this exposure depends on the judge model, the task, and the grading mode, so it must be measured for the specific setting the judge is deployed in, not assumed from a bias rate reported in a different study.


\section{The Agent Evaluation Problem}
\label{sec:agent-evaluation-problem}

An agent's task-solving process consists of a sequence of states, not a single output, which raises the question of what counts as the right unit of evaluation once the subject is an agent. The classical rational agent, introduced in Section~\ref{subsec:rationality}, offers a starting point for this question. Its behavior is evaluated by a performance measure, external to it and predetermined before it runs \cite[Chap.~2]{russellArtificialIntelligenceModern2022}.

\subsection{Same Outcome, Different Trajectories}
\label{subsec:quadratic-example}

Scoring the agent's behavior by its effect on the environment, rather than by the behavior itself, is the starting point for the example in Figures~\ref{fig:quadratic-tool} and \ref{fig:quadratic-direct}.

\begin{problembox}
  Two agents are given the equation $x^2 - 4 = 0$ and asked to report the values of $x$ that satisfy it.
\end{problembox}

\begin{scenariobox}{Scenario A}
  The agent calls a math tool, which solves the equation and returns the roots, as shown in Figure~\ref{fig:quadratic-tool}.

  \begin{figure}[H]
    \centering
    \begin{tikzpicture}[
        node distance=12mm,
        every node/.style={font=\small},
        box/.style={draw, rounded corners, align=center, minimum height=8mm, minimum width=20mm},
        arr/.style={-{Latex[length=2mm]}}
        ]
        \node[align=center, text width=22mm] (query) {Solve the equation\\$x^2-4=0$};
        \node[box, right=of query] (llm) {LLM};
        \node[align=center, text width=22mm, right=of llm] (output) {$x=2$, $x=-2$};
        \node[box, below=8mm of llm] (tool) {math-tool};
        \draw[arr] (query) -- (llm) node[midway, above] {\scriptsize 1};
        \draw[arr] ([xshift=-3mm]llm.south) -- ([xshift=-3mm]tool.north) node[midway, left] {\scriptsize 2};
        \draw[arr] ([xshift=3mm]tool.north) -- ([xshift=3mm]llm.south) node[midway, right] {\scriptsize 3};
        \draw[arr] (llm) -- (output) node[midway, above] {\scriptsize 4};
    \end{tikzpicture}
    \caption[Tool-call trajectory]{Tool-call trajectory.}
    \label{fig:quadratic-tool}
  \end{figure}
\end{scenariobox}

\begin{scenariobox}{Scenario B}
  The agent produces the roots directly, because the underlying language model has been trained on equations of this form and predicts the result without an intermediate step, as shown in Figure~\ref{fig:quadratic-direct}.

  \begin{figure}[H]
    \centering
    \begin{tikzpicture}[
        node distance=12mm,
        every node/.style={font=\small},
        box/.style={draw, rounded corners, align=center, minimum height=8mm, minimum width=20mm},
        unusedbox/.style={draw, dashed, rounded corners, align=center, minimum height=8mm, minimum width=20mm, text=gray, draw=gray},
        arr/.style={-{Latex[length=2mm]}}
        ]
        \node[align=center, text width=22mm] (query) {Solve the equation\\$x^2-4=0$};
        \node[box, right=of query] (llm) {LLM};
        \node[align=center, text width=22mm, right=of llm] (output) {$x=2$, $x=-2$};
        \node[unusedbox, below=8mm of llm] (tool) {math-tool};
        \draw[arr] (query) -- (llm) node[midway, above] {\scriptsize 1};
        \draw[arr] (llm) -- (output) node[midway, above] {\scriptsize 2};
    \end{tikzpicture}
    \caption[Direct-response trajectory]{Direct-response trajectory.}
    \label{fig:quadratic-direct}
  \end{figure}
\end{scenariobox}

Both scenarios report $x = 2$ and $x = -2$. An outcome-only criterion, applied to the final answer alone, scores them identically, because it has no access to the process that produced the answer and therefore cannot distinguish an answer produced by a tool call from one produced directly by the model. The two trajectories differ, however, in the guarantee they carry beyond this single equation. Scenario A routes equations to the same deterministic solver, so a harder or less common equation is handled by the same computation and solved with the same reliability. Scenario B carries no such guarantee: the model can produce incorrect roots for a harder or less common equation with the same confidence with which it reports the correct ones here.

\subsection{Outcome-Only Evaluation in Coding Benchmarks}
\label{subsec:coding-benchmarks-outcome-only}

Coding benchmarks that evaluate agents on software engineering tasks do not all treat the trajectory the same way: some score an agent purely by outcome, as in Scenario A and Scenario B above, while others also account for how that outcome was reached, and the choice determines what the benchmark measures. SWE-bench and ViBench, reviewed below, are outcome-only in their primary metric. SWE-bench scores an agent's submission by applying the proposed patch to the repository and running the associated test suite: a task instance is resolved when the patch applies and all tests pass, and the reported metric is the percentage of task instances resolved \cite{jimenezSWEbenchCanLanguage2024a}. This metric excludes the tool calls that the agent makes while producing the patch, their order, and the intermediate observations returned by the repository or the test runner. It reads only the final patch and the pass or fail outcome of the tests.

ViBench, a more recent coding benchmark, sharpens the same separation. Its primary score, Pass@1, credits an artifact only when every test plan associated with it scores full marks, indicating the absence of bugs or missing functionality \cite{zhongViBenchBenchmarkVibe2026a}. A test plan is a natural-language script of browser interactions that a human end user or QA tester would perform. Pass@1 verifies the running application's behavior through this browser automation, evaluated independently of the source code \cite{zhongViBenchBenchmarkVibe2026a}. The agent's tool-use trajectory is not part of this score: it is examined afterward, through what the benchmark calls \emph{trajectory-level} and \emph{qualitative failure-mode} analyses, presented separately from and after the Pass@1 and graded-score results \cite{zhongViBenchBenchmarkVibe2026a}. ViBench names this separation explicitly, treating trajectory analysis as a distinct, secondary step rather than a component folded into the metric. This separation may follow from a constraint on generic benchmark design. Trajectory analysis presupposes access to an agent's internal architecture. Because that architecture differs across the agents a benchmark evaluates, trajectory analysis cannot be standardized into a single metric the way an outcome check on the produced artifact can. This makes explicit the property that SWE-bench's design leaves implicit: the score reads only the outcome the agent produces. The process that produced it is analyzed, when it is analyzed at all, outside the scoring procedure.

\section{Trajectory-aware Evaluation}
\label{sec:trajectory-aware-evaluation}

Trajectory-aware evaluation scores the sequence of intermediate steps that an agent takes, its tool calls, plans, and decisions, rather than only the final output, and responds to the gap described in Section~\ref{sec:agent-evaluation-problem}.

The gap is not limited to trajectories that differ in whether a tool is invoked. It appears also between two agents that both use a tool: one agent can solve the equation with a single call to a symbolic math tool, while another can reach the same roots through multiple calls, applying a sequence of algebraic transformations and a general formula for the roots of a polynomial where the single call would have applied the specific formula for a quadratic. Both agents report $x = 2$ and $x = -2$, and an outcome-only criterion scores them identically. The two trajectories differ in the number of tool calls and in the computation that each call carries out, which affects the cost of the run. More generally, the correctness of the output is not sufficient to establish that the process producing it was carried out correctly. That has to be checked at the level of the trajectory itself.

Trajectory-aware evaluation is not a single method. Mohammadi et al.\ organized agent evaluation along two axes \cite{mohammadiEvaluationBenchmarkingLLM2025}:

\begin{itemize}
    \item \emph{What is evaluated}: outcome-oriented agent behavior, on one side, against process-oriented agent capabilities on the other, such as tool use, planning and reasoning, memory and context retention, and multi-agent collaboration.
    \item \emph{How the evaluation is carried out}: the interaction mode, the data, the metrics, and the tooling involved.
\end{itemize}

Within this taxonomy, the metrics and evaluation strategies that fit a given case depend on what is being assessed \cite{mohammadiEvaluationBenchmarkingLLM2025}. AgentBench, $\tau$-bench, and TRAJECT-Bench illustrate this variation with three concrete instances of trajectory-aware evaluation. All three target process-oriented capabilities, but they differ along the second axis of the taxonomy, the interaction mode used to elicit and score the trajectory, so each scores a different aspect of the agent's process:

\begin{itemize}
    \item \emph{AgentBench} evaluates agents across eight distinct environments, task settings such as an operating system shell, a database, and a knowledge graph, each with its own action space and metric \cite{liuAgentBenchEvaluatingLLMs2025}. The agent interacts with each environment over multiple turns, with the number of rounds ranging from five to thirty-five, so the benchmark assesses reasoning and decision-making across the interaction rather than at a single output.
    \item \emph{$\tau$-bench} measures an agent's ability to interact with a simulated user and with programmatic APIs while following a domain-specific policy document \cite{yao$t$benchBenchmarkToolAgentUser2024}. Its reward combines two terms: whether the database state, the content of the records the agent's tool calls wrote or changed, matches an annotated target state, and whether the agent's responses to the user adhere to the policy. The score thus depends on how the agent conducted the interaction, not only on the state it left behind.
    \item \emph{TRAJECT-Bench} reports trajectory-level diagnostics alongside final accuracy, including the correctness of tool selection and of the arguments passed to each tool, and whether dependencies and ordering constraints between steps are satisfied \cite{heTRAJECTBenchTrajectoryAwareBenchmark2025}.
\end{itemize}


\section{The Multi-Agent Trace}
\label{sec:multi-agent-trace}

The evaluation methods surveyed in Section~\ref{sec:trajectory-aware-evaluation} score a single agent's trajectory: its sequence of intermediate steps, tool calls, plans, and decisions. This framing does not cover what happens once several agents work together, as described in Chapter~\ref{chap:agents}. In a multi-agent system, agents exchange messages along a given topology. A single agent's trajectory records only that agent's own steps, so it captures neither the messages exchanged with other agents nor the topology they flow along. Evaluating such a system therefore requires a concept broader than a single trajectory, called the \emph{trace}, which consists of every participating agent's trajectory, the inter-agent messages exchanged between them, and the topology along which those messages and outputs flow. Trajectory-aware evaluation, as scoped in Section~\ref{sec:trajectory-aware-evaluation}, accounts for the first component of the trace; it does not account for the second and third.

Cemri et al. studied failures of multi-agent LLM systems using this broader concept, the execution trace, rather than individual agent trajectories \cite{cemriWhyMultiAgentLLM2025a}. Working from a set of traces produced by running multi-agent systems on their target tasks, they analyzed an initial set of 150 traces and derived from it a taxonomy of failure modes, the Multi-Agent System Failure Taxonomy (MAST) \cite{cemriWhyMultiAgentLLM2025a}. The taxonomy was then validated by annotating a larger set of 1642 traces against its failure modes, with inter-annotator agreement reported for the labeling procedure \cite{cemriWhyMultiAgentLLM2025a}. MAST groups 14 distinct failure modes into three categories \cite{cemriWhyMultiAgentLLM2025a}:

\begin{itemize}
    \item \emph{System Design Issues}: failures that originate from system design decisions, and from poor or ambiguous prompt specifications.
    \item \emph{Inter-Agent Misalignment}: failures that arise from a breakdown in the flow of information between agents during execution, such as a failure to share information, a misinterpretation of another agent's output, or a departure from an agreed task specification.
    \item \emph{Task Verification}: failures that involve inadequate verification of the final output, including no or incomplete verification, incorrect verification, or terminating the task before it is complete.
\end{itemize}

Inter-agent misalignment accounted for over 30\% of observed failures across the annotated set of 1642 traces \cite{cemriWhyMultiAgentLLM2025a}, and it is distinct from the other two categories, which are attributable to the design of a single agent or to the absence of a verification step. Since these failure modes are properties of the messages exchanged between agents and not of any single agent's trajectory, they are only visible when the trace is the object under evaluation.

\section{Failure Attribution}
\label{sec:failure-attribution}

The taxonomy in Section~\ref{sec:multi-agent-trace} names a failure mode but does not indicate which agent or which step in the trace caused it. Failure attribution is the additional task of assigning a trace failure to the agent and step responsible for it, so that a fix can target that agent or step rather than the run as a whole.

\subsection{Attribution Methods}
\label{sec:attribution-methods}

Zhang et al. formulated this step and introduced Who\&When, a benchmark of 127 annotated multi-agent systems in which each failed trace is labeled with the agent and the step responsible for the failure \cite{zhangWhichAgentCauses2025}. The automated methods that they evaluate on this benchmark are themselves built on the LLM-as-a-judge paradigm of Section~\ref{sec:llm-as-judge}, applied here to read a failure trace and return a verdict on who is responsible instead of grading a single output \cite{zhangWhichAgentCauses2025}. Against this benchmark, the best of these methods reached 53.5\% accuracy at identifying the responsible agent and 14.2\% accuracy at pinpointing the responsible step, and some methods performed below random on the step-level task \cite{zhangWhichAgentCauses2025}. The gap between the two figures indicates that identifying which agent is at fault is considerably easier than identifying at which point in that agent's trajectory the fault occurs. The low step-level accuracy is a further instance of the judge unreliability already established in Section~\ref{sec:llm-as-judge}, here applied to a harder task, reading a failure trace instead of grading a single response. These methods are trained and evaluated against the 127 traces in Who\&When. Labeling a trace at this granularity (agent and step) by hand is costly, which limits how much attribution data of this kind can exist.

AgenTracer targets this labeling cost: it automates the annotation of failed multi-agent trajectories through counterfactual replay and programmed fault injection, rather than relying on traces labeled by hand \cite{zhangAgenTracerWhoInducing2025}. Fault injection here means taking a trajectory that succeeded and deliberately corrupting one of its actions, so that the resulting failure has a known cause by construction \cite{zhangAgenTracerWhoInducing2025}. A tracer trained on the traces that the AgenTracer pipeline produces outperformed large proprietary judges, including Gemini-2.5-Pro and Claude-4-Sonnet, by up to 18.18\% on the Who\&When benchmark \cite{zhangAgenTracerWhoInducing2025}.

Both Who\&When's judges and AgenTracer's automated labels treat attribution as the selection of one responsible agent and one responsible step, a single ground truth to predict. DoVer questioned this framing and approached attribution as intervention-driven debugging instead: it generates a hypothesis about the responsible agent and step, then verifies the hypothesis through targeted interventions on the trace, in the same spirit as AgenTracer's fault injection but applied to check a hypothesis rather than to generate training data, such as editing a message or altering a plan, rather than accepting the hypothesis on its own \cite{maDoVerInterventionDrivenAuto2026}. The authors observed that framing attribution as the search for a single responsible agent and step is often ill-posed, since multiple distinct interventions can independently repair the same failed task \cite{maDoVerInterventionDrivenAuto2026}. In et al. reported the same difficulty from a different angle: existing benchmarks and methods largely assume a single deterministic root cause, but multi-agent system failures often admit multiple plausible attributions because of complex inter-agent dependencies and ambiguous execution trajectories \cite{inRethinkingFailureAttribution2026}. A trace failure, in this view, is not always reducible to one agent and one step; more than one attribution can be defensible for the same run.

\subsection{Cost of Attribution}
\label{sec:cost-of-attribution}

Attribution methods of the kind surveyed above need a complete trace to reason over, and reasoning over the trace once is often not enough: both AgenTracer and DoVer establish the responsible agent and step by replaying the trace under a modification rather than by reading it once. AgenTracer produces its labels through counterfactual replay combined with programmed fault injection \cite{zhangAgenTracerWhoInducing2025}, and DoVer replays each trial with an intervention applied in-situ, running three independent replays per intervention and, across the four settings it evaluates, between 1.4 and 4.0 replay trials per case \cite{maDoVerInterventionDrivenAuto2026}. Producing an attribution this way therefore implies two things. First, it requires re-executing part or all of the trace multiple times. Each replay checkpoints the state it starts from, and the procedure repeats until a hypothesis is confirmed. This adds a compute cost on top of whatever was already spent producing the original trace, and the cost scales with the number of interventions or fault injections a method tries before it converges on an attribution. Second, replaying a trace requires the trace to already exist in full: the run must have finished before any replay, checkpoint, or intervention can be performed on it. Attribution of this kind is therefore an offline procedure, carried out after the run is complete, and its cost in money and time is paid only after a failure has already happened. This implies a distinction between offline evaluation, performed on a completed trace, and online evaluation, performed during execution, developed in Section~\ref{sec:offline-online-evaluation}.


\section{Offline and Online Evaluation}
\label{sec:offline-online-evaluation}

\subsection{Definitions and Feedback Loop}
\label{sec:offline-online-definitions}

Xia et al. distinguished two stages at which a check can run, offline and online, and described evaluation-driven development and operations (EDDOps), an approach that places the two within a single closed feedback loop rather than treating them as alternatives \cite{xiaEvaluationDrivenDevelopmentOperations2025}. In EDDOps, offline and online evaluation results both feed a shared analysis step, which drives either a bounded runtime adjustment or a governed offline redevelopment. A redevelopment change is then re-verified with offline evaluation on the same cases before it reaches the live system \cite{xiaEvaluationDrivenDevelopmentOperations2025}.

\begin{itemize}
    \item \emph{Offline evaluation} runs against a predetermined suite of tasks or logged trajectories, on a subset of cases, ahead of release.
    \item \emph{Online evaluation} runs against live traffic, continuously, while the system operates in deployment.
\end{itemize}

A predetermined offline suite covers only a limited slice of the conditions a system meets after deployment. Xia et al. made the same point: predetermined suites cover a limited slice of possible tasks and inputs, so they are less likely to surface issues that emerge with changing data, tools, or user behavior \cite{xiaEvaluationDrivenDevelopmentOperations2025}. Online evaluation catches what the predetermined suite structurally cannot: conditions outside the suite, met only after deployment, such as emergent failures, drift, or violations of latency, cost, and availability targets \cite{xiaEvaluationDrivenDevelopmentOperations2025}.

\subsection{Security Evaluation}
\label{sec:security-evaluation}

Security is a further object of evaluation, and it motivates red-teaming, evaluation carried out by proactively attacking or testing a system to identify its vulnerabilities \cite{purpuraBuildingSafeGenAI2025}. Purpura et al. note that red-teaming is useful for an organization that wants to anticipate threats to its system and safeguard against them before production, but also that red-teaming alone does not guarantee the safety of a model after deployment, and recommend constant monitoring of a system once it is in production \cite{purpuraBuildingSafeGenAI2025}. Red-teaming therefore runs in both modes: an offline check performed once before release, and an online one run continuously against live traffic.

Purpura et al. distinguish red-teaming from guardrailing: red-teaming probes a system for vulnerabilities before deployment, and guardrailing safeguards the system after deployment \cite{purpuraBuildingSafeGenAI2025}. A guardrail is a check placed on a model's input or output, run online, that blocks a request or a response before it reaches the next step.

\subsection{Early Exit and the Feedback Loop}
\label{sec:early-exit-feedback-loop}

Running a check online, continuously, adds a capability an offline check cannot offer regardless of what it measures: the run can be stopped at the point the check fails, before the steps that would follow are executed. This is an alternative to the costly post-hoc attribution introduced in Section~\ref{sec:cost-of-attribution}: instead of reconstructing after the fact which step of a completed trace caused a failure, an online check identifies the failing step as it happens, and the run can exit there rather than run to completion. This early exit matters most for multi-agent systems and long trajectories. Without it, a run already known to have failed at one step would otherwise continue producing further steps, agent calls, and tool calls. The cost of those further steps is spent on a trajectory that cannot succeed. Stopping at the point of failure avoids this additional cost, and, since the state at that point is already checkpointed, the run can also be resumed or repaired from the failing step rather than discarded outright.

What happens after a check fails is decided by error handling, a step distinct from the check itself: the run is resumed from the checkpoint, repaired, or discarded, depending on the kind of failure that the check detected. Guardrailing is one instance of this pattern: it triggers the exit, and error handling determines the response that follows, such as refusing the request or returning a safe fallback.

This closes the loop described in Section~\ref{sec:offline-online-definitions}: evidence collected online about where and how a system fails is fed back to build the checks that later run online, so that a failure pattern identified once informs the check that catches and stops the same pattern the next time it occurs.


\section{Evaluating the Coordinator}
\label{sec:evaluating-coordinator}

The coordinator pattern of Section~\ref{sec:coordinator} fits the terms this chapter has built. The routing decision is a trajectory-level property in the sense of Section~\ref{sec:trajectory-aware-evaluation}: every sub-agent downstream of the coordinator acts on it, so a wrong routing decision is not repaired by anything that runs after it, regardless of how correctly that later step is carried out. The correctness of a run that involves a coordinator therefore cannot be read off the final output alone; it depends on whether the routing step itself was correct.

This also makes the routing decision an attribution target in the sense of Section~\ref{sec:failure-attribution}. Establishing that a run failed because of routing, rather than because of a later step, requires the same kind of attribution discussed there, with the same two limits: a judge cannot reliably confirm that a routing decision was wrong, and confirming it through replay requires the run to already be complete.

The offline and online regimes of Section~\ref{sec:offline-online-evaluation} both apply to the routing decision. Offline, routing decisions can be checked against a predetermined suite of requests with a known correct sub-agent. Online, a check on the routing decision offers a case of the early exit described in Section~\ref{sec:early-exit-feedback-loop}: a misrouting caught as the decision is made allows the run to be halted before any sub-agent acts on it, rather than allowing the run to complete on a request that was handed to the wrong sub-agent.

The experimental part of this work takes up the coordinator's routing decision under this framing.


